\documentclass[11pt]{article}

\usepackage[margin=1.35cm]{geometry}
\usepackage{amsmath,amssymb}
\usepackage{enumitem}
\usepackage{array}

\setlength{\columnsep}{0.65cm}
\setlength{\parindent}{0pt}
\setlist[enumerate]{leftmargin=*,itemsep=0.35em,topsep=0.25em}

% Multiple-choice bubble
\newcommand{\choice}[2]{%
    \par\noindent
    \(\bigcirc\)\ \textbf{#1.}\ #2
}

\newcommand{\bubblecell}[1]{\textbf{#1}\;\(\bigcirc\)A\;\(\bigcirc\)B\;\(\bigcirc\)C\;\(\bigcirc\)D}

\begin{document}

\begin{center}
{\Large\bfseries COSC 120: Introduction to Python Programming}\\[2pt]
{\large\bfseries In-Class Midterm Examination}\\[3pt]
Prof. Jordan C. Hanson
\end{center}

\vspace{0.1cm}

\noindent
\textbf{Name:} \hrulefill
\hspace{0.5cm}
\textbf{Student ID:} \fbox{\rule{0pt}{1.8ex}\hspace{4cm}}

\vspace{0.2cm}

\noindent
\textbf{Instructions:}
Choose the best answer for each question.
Fill in the circle next to exactly one answer, A--D.
Each question is worth 4 points.

\vspace{0.15cm}
\hrule
\vspace{0.2cm}

\begin{center}
{\large\bfseries Answer Sheet}
\end{center}

\small
\renewcommand{\arraystretch}{1.3}
\begin{center}
\begin{tabular}{|m{0.18\textwidth}|m{0.18\textwidth}|m{0.18\textwidth}|m{0.18\textwidth}|m{0.18\textwidth}|}
\hline
\bubblecell{1} & \bubblecell{2} & \bubblecell{3} & \bubblecell{4} & \bubblecell{5} \\ \hline
\bubblecell{6} & \bubblecell{7} & \bubblecell{8} & \bubblecell{9} & \bubblecell{10} \\ \hline
\bubblecell{11} & \bubblecell{12} & \bubblecell{13} & \bubblecell{14} & \bubblecell{15} \\ \hline
\bubblecell{16} & \bubblecell{17} & \bubblecell{18} & \bubblecell{19} & \bubblecell{20} \\ \hline
\bubblecell{21} & \bubblecell{22} & \bubblecell{23} & \bubblecell{24} & \bubblecell{25} \\ \hline
\end{tabular}
\end{center}
\normalsize

\vspace{0.15cm}
\hrule
\vspace{0.25cm}

\twocolumn

%=========================================================
\section*{Python Fundamentals}

\begin{enumerate}

\item What is printed by the following code?

\begin{verbatim}
x = 7
y = 3
print(x // y, x % y)
\end{verbatim}

\choice{A}{2 1}
\choice{B}{2 3}
\choice{C}{2.33 1}
\choice{D}{3 1}

\item What is printed?

\begin{verbatim}
x = 3.14159
print(f"{x:.2f}")
\end{verbatim}

\choice{A}{3.1}
\choice{B}{3.14}
\choice{C}{3.142}
\choice{D}{3.14159}

\item Consider

\begin{verbatim}
student = ("Maya", 20681001, 5)
\end{verbatim}

What is the value of

\begin{verbatim}
student[2]
\end{verbatim}

\choice{A}{Maya}
\choice{B}{20681001}
\choice{C}{5}
\choice{D}{An error occurs}

\item What is the final value of \texttt{numbers}?

\begin{verbatim}
numbers = [1, 2]
numbers.append(3)
\end{verbatim}

\choice{A}{[1, 2]}
\choice{B}{[1, 2, 3]}
\choice{C}{[3, 1, 2]}
\choice{D}{1, 2, 3}

\item What decimal value is produced by

\begin{verbatim}
int("1F", 16)
\end{verbatim}

\choice{A}{15}
\choice{B}{16}
\choice{C}{31}
\choice{D}{32}

\end{enumerate}

%=========================================================
\section*{Boolean Logic and Decisions}

\begin{enumerate}
\setcounter{enumi}{5}

\item Consider the placement code

\begin{verbatim}
score = 3

if score >= 4:
    course = "Calculus"
elif score >= 2:
    course = "Trigonometry"
else:
    course = "Algebra"
\end{verbatim}

What is stored in \texttt{course}?

\choice{A}{Calculus}
\choice{B}{Trigonometry}
\choice{C}{Algebra}
\choice{D}{Nothing}

\item What is the value of

\begin{verbatim}
True or False and False
\end{verbatim}

Recall that \texttt{and} is evaluated before \texttt{or}.

\choice{A}{True}
\choice{B}{False}
\choice{C}{0}
\choice{D}{An error occurs}

\item A screening function contains

\begin{verbatim}
age > 70 and
weight > 200 and
gender == "M"
\end{verbatim}

Which patient satisfies the condition?

\choice{A}{Male, age 70, weight 210}
\choice{B}{Male, age 75, weight 195}
\choice{C}{Female, age 75, weight 220}
\choice{D}{Male, age 75, weight 220}

\end{enumerate}

%=========================================================
\section*{Loops and Nested Lists}

\begin{enumerate}
\setcounter{enumi}{8}

\item What values are printed?

\begin{verbatim}
for i in range(1, 4):
    print(i)
\end{verbatim}

\choice{A}{0, 1, 2, 3}
\choice{B}{1, 2, 3}
\choice{C}{1, 2, 3, 4}
\choice{D}{0, 1, 2}

\item How many times does this loop execute?

\begin{verbatim}
i = 0

while i < 3:
    print(i)
    i = i + 1
\end{verbatim}

\choice{A}{2}
\choice{B}{3}
\choice{C}{4}
\choice{D}{The loop never stops}

\item Consider

\begin{verbatim}
for i in range(len(courses)):

    for j in range(i + 1,
                   len(courses)):
        ...
\end{verbatim}

Why does the inner loop begin at \texttt{i + 1}?

\choice{A}{To skip every other course}
\choice{B}{To compare each course only with itself}
\choice{C}{To avoid comparing a pair twice and avoid comparing a course with itself}
\choice{D}{Because Python lists begin at index 1}

\item Consider the matrix

\begin{verbatim}
A = [
 [4, 8, 2],
 [7, 5, 9],
 [1, 6, 3]
]
\end{verbatim}

What is the value of

\begin{verbatim}
A[1][2]
\end{verbatim}

\choice{A}{5}
\choice{B}{7}
\choice{C}{9}
\choice{D}{6}

\end{enumerate}

%=========================================================
\section*{Functions}

\begin{enumerate}
\setcounter{enumi}{12}

\item Consider

\begin{verbatim}
def square(x):
    return x * x

square(5)
\end{verbatim}

Which statement is correct?

\choice{A}{\texttt{x} is the argument and 5 is the parameter}
\choice{B}{\texttt{x} is the parameter and 5 is the argument}
\choice{C}{Both \texttt{x} and 5 are parameters}
\choice{D}{Neither is a parameter}

\item Why is \texttt{return} useful in a function such as

\begin{verbatim}
def find_patient(...):
    ...
    return patient
\end{verbatim}

\choice{A}{It permanently prints the patient}
\choice{B}{It sends a value back to the calling code}
\choice{C}{It automatically creates a loop}
\choice{D}{It imports another module}

\item In our Fourier-series exercise, the functions were
organized as

\begin{verbatim}
plot_series()
sawtooth()
fourier_term()
\end{verbatim}

Which call sequence correctly describes the program?

\choice{A}{\texttt{fourier\_term()} $\rightarrow$ \texttt{sawtooth()} $\rightarrow$ \texttt{plot\_series()}}
\choice{B}{\texttt{plot\_series()} $\rightarrow$ \texttt{sawtooth()} $\rightarrow$ \texttt{fourier\_term()}}
\choice{C}{\texttt{sawtooth()} $\rightarrow$ \texttt{plot\_series()} $\rightarrow$ \texttt{fourier\_term()}}
\choice{D}{None of the functions call one another}

\end{enumerate}

%=========================================================
\section*{Modules}

\begin{enumerate}
\setcounter{enumi}{15}

\item Suppose \texttt{math\_tools.py} contains a function called
\texttt{square()}. Which code correctly imports the module and
calls the function?

\choice{A}{\texttt{import math\_tools}\\ \texttt{math\_tools.square(5)}}
\choice{B}{\texttt{import square}\\ \texttt{math\_tools(5)}}
\choice{C}{\texttt{open math\_tools}\\ \texttt{square.math\_tools(5)}}
\choice{D}{\texttt{module math\_tools}\\ \texttt{square(5)}}

\item Which code correctly imports only the function
\texttt{square()}?

\choice{A}{\texttt{import square from math\_tools}}
\choice{B}{\texttt{from math\_tools import square}}
\choice{C}{\texttt{from square import math\_tools}}
\choice{D}{\texttt{math\_tools import(square)}}

\item Before a custom file such as

\begin{verbatim}
course_tools.py
\end{verbatim}

can be imported into a Google Colab session, what must normally
be done?

\choice{A}{Convert it to a PDF}
\choice{B}{Upload the Python file to the Colab runtime}
\choice{C}{Rename it with a \texttt{.txt} extension}
\choice{D}{Place all functions inside \texttt{print()}}

\end{enumerate}

%=========================================================
\section*{Strings and Searching}

\begin{enumerate}
\setcounter{enumi}{18}

\item What is the value of

\begin{verbatim}
"Garcia" in "Maria Garcia"
\end{verbatim}

\choice{A}{True}
\choice{B}{False}
\choice{C}{Garcia}
\choice{D}{An error occurs}

\item Which condition performs a case-insensitive partial-name
search?

\choice{A}{\texttt{search == patient[1]}}
\choice{B}{\texttt{search in patient[1]}}
\choice{C}{\texttt{search.lower() in patient[1].lower()}}
\choice{D}{\texttt{search.upper() == patient[1].lower()}}

\item A function searches every patient record but finds no
matching patient ID. What value did our exercise specify should
be returned?

\choice{A}{0}
\choice{B}{False}
\choice{C}{""}
\choice{D}{None}

\end{enumerate}

%=========================================================
\section*{Course Lists and Sorting}

\begin{enumerate}
\setcounter{enumi}{21}

\item A course is stored as

\begin{verbatim}
course = [
 "PHYS", 100, "SLC", 202,
 "MWF", 1400
]
\end{verbatim}

What is the value of \texttt{course[5]}?

\choice{A}{SLC}
\choice{B}{202}
\choice{C}{MWF}
\choice{D}{1400}

\item Suppose

\begin{verbatim}
def get_time(course):
    return course[5]

def sort_by_time(courses):
    courses.sort(key=get_time)
\end{verbatim}

What does \texttt{key=get\_time} tell Python to do?

\choice{A}{Sort using the course meeting time}
\choice{B}{Delete the meeting time}
\choice{C}{Sort only courses named ``Time''}
\choice{D}{Replace each course with its meeting time}

\item Which time would appear first after sorting numerically?

\begin{verbatim}
930, 1400, 1100, 1530
\end{verbatim}

\choice{A}{930}
\choice{B}{1100}
\choice{C}{1400}
\choice{D}{1530}

\item Consider these two courses:

\begin{verbatim}
["PHYS",100,"SLC",202,
 "MWF",1400]

["MATH",200,"SLC",202,
 "TR",1400]
\end{verbatim}

Using the room-conflict rule from our course-schedule exercise,
do these courses conflict?

\choice{A}{Yes, because the room is the same}
\choice{B}{Yes, because the time is the same}
\choice{C}{No, because they do not share a meeting day}
\choice{D}{No, because the course numbers are different}

\end{enumerate}

\vfill

\begin{center}
\textbf{End of Exam}

Check that exactly one answer is clearly marked for each question.
\end{center}

\end{document}